\chapter{The Structure of a Thesis in Computer Science}\label{structure}

A thesis in computer science follows the structure of academic writing and consists of several chapters, each chapter contains several sections, and each section can contain several subsections. Sections and subsections contain several paragraphs of text as well as, for example, lists, tables, and figures. It is recommended to start your thesis by studying the scientific toolbox~\footnote{Find our notes on scientific work at \url{https://webis.de/lecturenotes.html\#part-scientific-toolbox}}.

\enlargethispage{\baselineskip}

\section{References}\label{references}

Elements like sections, figures, or tables can be referenced (see Section~\ref{references}) by assigning a \texttt{\textbackslash label\{label-name\}} after the \texttt{\textbackslash section} command or within a \texttt{\textbackslash begin\{\dots\}} environment and referencing the label with \texttt{\textbackslash ref\{label-name\}}.

\section{Links and Citations}

In scientific work, all external sources must be cited or linked.

\paragraph{Scientific Articles} Papers, Textbooks, theses, or other scientific work should always be cited as \texttt{@article}, \texttt{@inproceedings}, or \texttt{@book} via biblatex. To cite a work, add a corresponding bibliography entry to your \texttt{MasterThesisGuardia.bib} file and reference it in the text.

Using \texttt{biblatex}, you can display citations in multiple ways depending on your sentence structure:
\begin{itemize}
    \item Standard citation with \texttt{\textbackslash cite\{bib-key\}}: \cite{kabanov51Attacks602026}
    \item Parenthetical citation with \texttt{\textbackslash parencite\{bib-key\}}: \parencite{manning:2001}
    \item Narrative/Textual citation with \texttt{\textbackslash textcite\{bib-key\}}: \textcite{manning:2001}
    \item Author-only extraction with \texttt{\textbackslash citeauthor\{bib-key\}}: \citeauthor{manning:2001}
\end{itemize}

\paragraph{Non-scientific Articles} Journalistic articles, books, blog posts, and various web sources with a known author and title should generally be cataloged using the \texttt{@online} or \texttt{@misc} entry types in your bibliography database. For web sources, make sure to populate the \texttt{url} field and provide the exact retrieval date using the \texttt{urldate = {YYYY-MM-DD}} field format so it displays cleanly under French localization rules.

\paragraph{Other Sources} Other secondary assets, like standalone images, external libraries, or specific software source code deployments, can be dynamically referenced by providing an explicit link and a date of last access within a footnote on the current page.